\chapter{Thalassaemia}
\index{Thalassaemia}

\medskip
\noindent\textit{Educational reference; always seek professional advice for individual care.}

\section*{Definition and Overview}
Thalassaemia is a group of inherited blood disorders in which the body makes abnormal or reduced amounts of haemoglobin, the oxygen-carrying protein in red blood cells. Faulty genes for the alpha or beta globin chains lead to ineffective red cell production, chronic anaemia, and a variable burden of complications. Severity ranges from silent carrier states and mild thalassaemia trait (thalassaemia minor), which often cause little more than a small red cell size on blood tests, through thalassaemia intermedia, to transfusion-dependent thalassaemia major, which requires lifelong specialist care. Alpha and beta thalassaemia are the main categories. In many countries, thalassaemia is more common in people with ancestry from the Mediterranean, Middle East, Indian subcontinent, Southeast Asia, Africa, and Pacific regions, reflecting historical malaria selection. Carrier screening before or early in pregnancy, accurate diagnosis, and modern transfusion and iron-chelation programmes have transformed outcomes. This chapter outlines types, inheritance, clinical features, and care pathways for patients and families.

\section*{Epidemiology}
Globally, millions carry thalassaemia genes, and hundreds of thousands live with severe disease. In many countries, trait is relatively common in certain ethnic communities; severe disease is less common but well represented in paediatric and adult haematology services. Beta thalassaemia major historically affected Mediterranean-origin families and now also many Asian-origin families. Alpha thalassaemia deletions are frequent in Southeast Asian and some Pacific and African populations; the most severe form, haemoglobin Bart's hydrops fetalis (no functional alpha genes), is usually fatal in utero or soon after birth without highly specialised intervention. Partner testing when one person is a carrier is standard public health practice. Migration has broadened the genetic mix seen in local clinics, so ethnicity alone must not exclude screening when microcytosis is found.

\section*{Aetiology and Pathophysiology}
Haemoglobin A in adults is mostly two alpha and two beta chains. Reduced synthesis of one chain type leaves an excess of the other, damaging red cell precursors in the bone marrow and shortening red cell survival.

\begin{itemize}
    \item \textbf{Beta thalassaemia:} mutations reduce or abolish beta globin; excess alpha chains precipitate, causing ineffective erythropoiesis and haemolysis; severity depends on residual beta output (beta-plus versus beta-zero) and co-inheritance of other haemoglobin variants such as HbE
    \item \textbf{Alpha thalassaemia:} usually gene deletions; one or two gene loss causes silent carrier or trait; three-gene loss causes haemoglobin H disease; four-gene loss causes Bart's hydrops
    \item \textbf{Ineffective erythropoiesis:} marrow expansion tries to compensate, which can widen bones and increase metabolic demand
    \item \textbf{Chronic anaemia and hypoxia:} drive increased iron absorption even before transfusions
    \item \textbf{Transfusional iron overload:} regular red cell transfusions deposit iron in liver, heart, and endocrine glands unless chelated
    \item \textbf{Inheritance:} autosomal patterns; two carrier parents have a one-in-four chance of a child with major disease each pregnancy (for classic recessive pairings)
\end{itemize}

\section*{Clinical Features}
\begin{itemize}
    \item \textbf{Thalassaemia trait:} often asymptomatic; mild anaemia or none; small pale red cells; sometimes fatigue misattributed to iron deficiency
    \item \textbf{Thalassaemia intermedia / non-transfusion-dependent:} moderate anaemia, intermittent transfusion need, leg ulcers, bone changes, gallstones, pulmonary hypertension risk, and iron loading from absorption
    \item \textbf{Beta thalassaemia major (untreated historical picture):} severe anaemia from infancy, failure to thrive, frontal bossing, maxillary prominence, enlarged liver and spleen
    \item \textbf{Transfusion-dependent disease on modern care:} children grow much better but need scheduled transfusions and chelation; breakthrough symptoms if delayed
    \item \textbf{Iron overload signs:} liver disease, cardiomyopathy, diabetes, hypogonadism, hypothyroidism, bronze skin pigmentation if poorly chelated
    \item \textbf{Splenomegaly and hypersplenism:} low platelets and white cells, increased transfusion requirement
    \item \textbf{Haemoglobin H disease:} variable anaemia, worse with infection or oxidant stress; splenomegaly common
\end{itemize}

\section*{Differential Diagnosis}
\begin{itemize}
    \item \textbf{Iron deficiency anaemia:} also microcytic; distinguished by iron studies, response to iron, and haemoglobin electrophoresis or DNA testing---giving iron long-term in thalassaemia trait without deficiency is unhelpful and potentially harmful
    \item \textbf{Anaemia of chronic disease:} usually normocytic or mildly microcytic with different iron pattern
    \item \textbf{Sideroblastic anaemia and lead poisoning:} microcytosis with other clues
    \item \textbf{Other haemoglobinopathies:} sickle cell disease, HbC, HbE disease---often coexist or compound with thalassaemia alleles
    \item \textbf{Hereditary spherocytosis and other haemolytic anaemias:} different red cell indices and film
\end{itemize}

\section*{When to Seek Medical Care}
Seek GP or haematology advice for unexplained microcytic anaemia, family history of thalassaemia, or preconception planning. Pregnant people with trait need partner testing early. Known patients need urgent review for fever (asplenia risk if splenectomised), severe anaemia symptoms, chest pain, or heart failure signs.

\textbf{Contact your local emergency services for:}
\begin{itemize}
    \item Severe shortness of breath, chest pain, collapse, or signs of heart failure
    \item High fever after splenectomy or in a severely anaemic child
    \item Stroke-like symptoms or severe abdominal pain with known haemoglobinopathy complications
\end{itemize}

\section*{Diagnosis and Investigations}
\begin{itemize}
    \item \textbf{Full blood count and film:} low MCV and MCH out of proportion to anaemia in trait; target cells, basophilic stippling, and nucleated red cells in severe disease
    \item \textbf{Iron studies:} to avoid misdiagnosing trait as iron deficiency and to monitor overload
    \item \textbf{Haemoglobin analysis:} HPLC or electrophoresis shows raised HbA2 in classic beta trait; HbH inclusions in alpha disease; patterns vary with age and transfusion
    \item \textbf{DNA testing:} definitive alpha genotyping and complex beta mutations; essential for prenatal diagnosis
    \item \textbf{Partner and cascade testing:} family screening after a carrier is found
    \item \textbf{Prenatal and preimplantation options:} CVS or amniocentesis genetics; IVF with PGT in some couples---counselling is specialised
    \item \textbf{Monitoring severe disease:} ferritin, MRI for liver and cardiac iron, echocardiogram, endocrine tests, bone density, viral serology related to transfusion history
\end{itemize}

\section*{Management}
\begin{itemize}
    \item \textbf{Trait:} usually no treatment; genetic counselling; treat true iron deficiency only if proven; avoid unnecessary iron
    \item \textbf{Regular transfusion programmes:} aim to suppress ineffective erythropoiesis and support growth and activity in major disease
    \item \textbf{Iron chelation:} desferrioxamine, deferasirox, or deferiprone regimens tailored to iron burden and tolerability
    \item \textbf{Folic acid:} often supplemented because of high red cell turnover
    \item \textbf{Splenectomy:} selected cases with hypersplenism; requires vaccines and lifelong infection-awareness
    \item \textbf{Hydroxyurea and other disease-modifiers:} used in some non-transfusion-dependent or selected patients
    \item \textbf{Haematopoietic stem cell transplantation:} potentially curative in eligible children with major disease and suitable donors
    \item \textbf{Emerging gene therapy:} available in limited specialised settings internationally; eligibility and access evolve
    \item \textbf{Supportive care:} endocrine replacement, cardiac care, bone health, psychosocial support, and transition to adult services
\end{itemize}

\section*{Complications}
\begin{itemize}
    \item \textbf{Iron overload organ damage:} cardiomyopathy, cirrhosis, diabetes, infertility, hypothyroidism, hypoparathyroidism
    \item \textbf{Transfusion reactions and alloimmunisation:} difficulty matching blood
    \item \textbf{Infections:} historically hepatitis viruses; bacterial risk after splenectomy
    \item \textbf{Bone disease:} osteoporosis, fractures, residual facial bone changes if late treatment
    \item \textbf{Thrombosis and pulmonary hypertension:} especially after splenectomy or in intermedia
    \item \textbf{Leg ulcers and extramedullary masses:} in inadequately suppressed erythropoiesis
\end{itemize}

\section*{Prognosis and Outcomes}
With modern transfusion and chelation, many people with beta thalassaemia major live into middle age and beyond with good quality of life, education, and employment. Prognosis hinges on adherence to chelation and access to specialist centres. Trait carriers have a normal life expectancy. Haemoglobin H disease varies from mild to transfusion-needing. Bart's hydrops remains a high-risk obstetric and neonatal emergency. Genetic counselling enables informed reproductive choices.

\section*{Prevention and Self-Care}
\begin{itemize}
    \item Accept carrier screening if you have microcytosis, relevant ancestry, or a partner who is a carrier
    \item Plan pregnancy with known carrier status so partner testing and prenatal options are timely
    \item Attend all transfusion and chelation appointments; report side effects early
    \item Keep vaccinations current, especially if splenectomised; seek urgent care for fever
    \item Avoid unsupervised iron supplements
    \item Maintain bone health with diet, vitamin D as advised, weight-bearing activity, and no smoking
    \item Carry a medical summary listing diagnosis, antibodies, and chelation drugs when travelling
\end{itemize}

\section*{Sources and Further Reading}
\begin{itemize}
    \item NHS. \textit{Thalassaemia.} https://www.nhs.uk/
    \item Mayo Clinic. \textit{Thalassemia.} https://www.mayoclinic.org/
    \item MSD Manual. \textit{Thalassemias.} https://www.msdmanuals.com/
    \item MedlinePlus. \textit{Thalassemia.} https://medlineplus.gov/
    \item Better Health Channel. \textit{Thalassaemia.} https://www.betterhealth.vic.gov.au/
    \item Healthdirect. \textit{Thalassaemia.} https://www.healthdirect.gov.au/
\end{itemize}
